\documentclass[10pt, letterpaper]{article}

% Packages:
\usepackage[
    ignoreheadfoot,
    top=0.9 cm,
    bottom=0.9 cm,
    left=1.3 cm,
    right=1.3 cm,
    footskip=0.6 cm,
]{geometry}
\usepackage[explicit]{titlesec}
\usepackage{tabularx}
\usepackage{array}
\usepackage[dvipsnames]{xcolor}
\definecolor{primaryColor}{RGB}{217, 119, 87}
\usepackage{enumitem}
\usepackage{fontawesome5}
\usepackage{amsmath}
\usepackage[
    pdftitle={Andrés Ortega CV},
    pdfauthor={Andrés González Ortega},
    pdfcreator={LaTeX with RenderCV},
    colorlinks=true,
    urlcolor=primaryColor
]{hyperref}
\usepackage[pscoord]{eso-pic}
\usepackage{calc}
\usepackage{bookmark}
\usepackage{lastpage}
\usepackage{changepage}
\usepackage{paracol}
\usepackage{ifthen}
\usepackage{needspace}
\usepackage{iftex}

% ATS parsable:
\ifPDFTeX
    \input{glyphtounicode}
    \pdfgentounicode=1
    \usepackage[T1]{fontenc}
    \usepackage[utf8]{inputenc}
    \usepackage{lmodern}
\fi

\usepackage[default, type1]{sourcesanspro}

% Settings:
\AtBeginEnvironment{adjustwidth}{\partopsep0pt}
\pagestyle{empty}
\setcounter{secnumdepth}{0}
\setlength{\parindent}{0pt}
\setlength{\topskip}{0pt}
\setlength{\columnsep}{0.15cm}
\makeatletter
\let\ps@customFooterStyle\ps@plain
\patchcmd{\ps@customFooterStyle}{\thepage}{
    \color{gray}\textit{\small Andrés Ortega -- Página \thepage{} de \pageref*{LastPage}}
}{}{}
\makeatother
\pagestyle{customFooterStyle}

\titleformat{\section}{
    \needspace{4\baselineskip}
    \large\color{primaryColor}
}{
}{
}{
    \textbf{#1}\hspace{0.15cm}\titlerule[0.8pt]\hspace{-0.1cm}
}[]

\titlespacing{\section}{
    -1pt
}{
    0.06 cm
}{
    0.04 cm
}

\newenvironment{highlights}{
    \begin{itemize}[
        topsep=0.02 cm,
        parsep=0.05 cm,
        partopsep=0pt,
        itemsep=0pt,
        leftmargin=0.4 cm + 10pt
    ]
}{
    \end{itemize}
}

\newenvironment{highlightsforbulletentries}{
    \begin{itemize}[
        topsep=0.02 cm,
        parsep=0.02 cm,
        partopsep=0pt,
        itemsep=0pt,
        leftmargin=10pt
    ]
}{
    \end{itemize}
}

\newenvironment{onecolentry}{
    \begin{adjustwidth}{
        0.2 cm + 0.00001 cm
    }{
        0.2 cm + 0.00001 cm
    }
}{
    \end{adjustwidth}
}

\newenvironment{twocolentry}[2][]{
    \onecolentry
    \def\secondColumn{#2}
    \setcolumnwidth{\fill, 4.0 cm}
    \begin{paracol}{2}
}{
    \switchcolumn \raggedleft \secondColumn
    \end{paracol}
    \endonecolentry
}

\newenvironment{twocolentrynarrow}[2][]{
    \onecolentry
    \def\secondColumn{#2}
    \setcolumnwidth{\fill, 2.2 cm}
    \begin{paracol}{2}
}{
    \switchcolumn \raggedleft \secondColumn
    \end{paracol}
    \endonecolentry
}

\newenvironment{threecolentry}[3][]{
    \onecolentry
    \def\thirdColumn{#3}
    \setcolumnwidth{1 cm, \fill, 4.5 cm}
    \begin{paracol}{3}
    {\raggedright #2} \switchcolumn
}{
    \switchcolumn \raggedleft \thirdColumn
    \end{paracol}
    \endonecolentry
}

\newenvironment{header}{
    \setlength{\topsep}{0pt}\par\kern\topsep\centering\color{primaryColor}\linespread{1.5}
}{
    \par\kern\topsep
}

% QR code + date in top-right corner
\newcommand{\placetopright}{%
  \AddToShipoutPictureFG*{%
    \put(\LenToUnit{\paperwidth-0.3 cm},\LenToUnit{\paperheight-0.2 cm}){%
      \vtop{{\null}\makebox[0pt][r]{%
        \vbox{%
          \hbox to 1.3cm{\hfil{\small\color{gray}\textit{Abril 2026}}\hfil}%
          \vspace{0.05cm}%
          \hbox{\pdfximage width 1.3cm {../qr_portfolio.png}\pdfrefximage\pdflastximage}%
          \vspace{0.03cm}%
          \hbox to 1.3cm{\hfil{\tiny\color{gray}\textit{portafolio}}\hfil}%
        }
      }}%
    }%
  }%
}

% Colored highlight for key terms
\newcommand{\hl}[1]{\textcolor{primaryColor}{\textbf{#1}}}

\let\hrefWithoutArrow\href
\renewcommand{\href}[2]{\hrefWithoutArrow{#1}{\ifthenelse{\equal{#2}{}}{ }{#2 }\raisebox{.15ex}{\footnotesize \faExternalLink*}}}


\begin{document}
    \newcommand{\AND}{\unskip
        \cleaders\copy\ANDbox\hskip\wd\ANDbox
        \ignorespaces
    }
    \newsavebox\ANDbox
    \sbox\ANDbox{}

    \placetopright

    \begin{header}
        \fontsize{24 pt}{24 pt}
        \textbf{Andrés González Ortega}

        \vspace{0.15 cm}

        \normalsize
        \mbox{{\footnotesize\faMapMarker*}\hspace*{0.13cm}CDMX}%
        \kern 0.25 cm%
        \AND%
        \kern 0.25 cm%
        \mbox{\hrefWithoutArrow{mailto:gonorandres@ciencias.unam.mx}{{\footnotesize\faEnvelope[regular]}\hspace*{0.13cm}gonorandres@ciencias.unam.mx}}%
        \kern 0.25 cm%
        \AND%
        \kern 0.25 cm%
        \mbox{\hrefWithoutArrow{tel:+525548344672}{{\footnotesize\faPhone*}\hspace*{0.13cm}+52 55 4834 4672}}%
        \AND%
        \kern 0.25 cm%
        \mbox{\hrefWithoutArrow{https://gonorandres.github.io/\#proyectos}{{\footnotesize\faLink}\hspace*{0.13cm}Mi web (25+ proyectos)}}%
        \kern 0.25 cm%
        \AND%
        \kern 0.25 cm%
        \mbox{\hrefWithoutArrow{https://github.com/GonorAndres}{{\footnotesize\faGithub}\hspace*{0.13cm}GitHub}}%
    \end{header}

    \vspace{0.03 cm}

    \section{Sobre mí}

        \begin{onecolentry}
Soy actuario de la UNAM con \textbf{SOA Exam Probability aprobado}, becario de ciencia de datos en \textbf{Grupo Gigante} y co-autor de artículo científico en el IGF-UNAM. Tomo problemas de seguros, banca y retail (tarificación, reservas, scoring crediticio, segmentación) y los llevo a producción: modelos \textbf{estadísticos} y de \textbf{machine learning}, pipelines en \textbf{GCP}, APIs y dashboards. Practico la formación continua (preparo la certificación \textbf{Databricks Machine Learning Engineer Associate}) y disfruto el trabajo colaborativo con equipos multidisciplinarios. En \href{https://gonorandres.github.io/blog/}{mi página web} comparto lo que aprendo y construyo.
        \end{onecolentry}


    \section{Experiencia}

        \begin{onecolentry}
            \textbf{Becario Ciencia de Datos} -- Grupo Gigante\quad{\small\color{gray}\textit{Septiembre 2025 -- Presente}}
            \begin{highlights}
                \item Trabajo \textbf{hands-on} en el \textbf{CDP} del grupo con \textbf{Python} y \textbf{GCP}: resolución de identidad entre marcas, entrenamiento de modelos de \textbf{clasificación} para anticipar abandono y construcción de pipelines \textbf{ETL} (BigQuery, Cloud SQL, Firebase) sobre millones de perfiles. Proyectos transversales: análisis de inventarios, comportamiento web y activación de audiencias.
            \end{highlights}
            \end{onecolentry}

        \vspace{0.02 cm}

        \begin{onecolentry}
            \textbf{Asistente de Investigación Científica} -- Instituto de Geofísica, UNAM \textit{(Servicio Social)}\quad{\small\color{gray}\textit{2025 -- Presente}}
            \begin{highlights}
                \item Co-autor de artículo científico sobre predicción de erupciones del Popocatépetl con \textbf{procesos de renovación Weibull}: modelado estadístico de eventos raros de alto impacto. Supervisión: \hrefWithoutArrow{https://scholar.google.com/citations?user=lDD0DZQAAAAJ}{Dr.\ Hugo Delgado-Granados} (IGF-UNAM).
            \end{highlights}
            \end{onecolentry}


    \section{Habilidades}

\begin{onecolentry}
    \begin{highlightsforbulletentries}
        \item \textit{Lenguajes:} Python, R, SQL, TypeScript, Bash, Git, \LaTeX, Excel Avanzado, Power BI
        \item \textit{Ciencia Actuarial:} \textbf{Tarificación} (GLM, frecuencia-severidad, credibilidad), \textbf{Reservas} (BEL, IBNR, Chain Ladder, BF), Lee-Carter, \textbf{RCS/SCR} bajo \textbf{LISF/CUSF}, EMSSA-09, Reaseguro, CONAPO/HMD, Pensiones IMSS (Ley 73/97, Fondo Bienestar)
        \item \textit{Estadística \& Ciencia de Datos:} scikit-learn, PyTorch, XGBoost, LightGBM, GLM Tweedie, Monte Carlo, SHAP, A/B testing, inferencia Bayesiana, Pandas, FastAPI, Streamlit, Plotly, Transformers
        \item \textit{Data Engineering \& Cloud:} GCP (BigQuery, Pub/Sub, Cloud Run), Dagster, Apache Beam, Databricks, PySpark, Neo4j, PostgreSQL, Docker, Terraform
        \item \textit{IA \& Agentes:} LLMs, Claude API, Claude Agent SDK, RAG, FTS5/BM25
        \item \textit{Idiomas:} Inglés avanzado (lectura, redacción, comunicación fluida)
    \end{highlightsforbulletentries}
\end{onecolentry}


    \section{Educación}

        \begin{twocolentry}{2021 -- 2025}
            \textbf{UNAM, Facultad de Ciencias} -- Licenciatura en Actuaría. Pasante (100\% de créditos cubiertos); titulación por exámenes profesionales SOA (\href{https://www.soa.org/}{Society of Actuaries}).
        \end{twocolentry}

    \section{Certificaciones}

    \begin{onecolentry}
        \begin{highlightsforbulletentries}
            \item \textbf{SOA Exam Probability} -- Aprobado Marzo 2026 \href{https://drive.google.com/file/d/1rt3emgBnPpQi7NiXkBQeA5WwJTV0JuAf/view?usp=sharing}{}
            \item \textbf{Associate Data Analyst} -- DataCamp \href{https://drive.google.com/drive/folders/1-QUuVH3zPBPsl6RhUvuUWxpZkP1i8TYb?usp=sharing}{}
            \item \textbf{Associate Data Scientist in R} -- DataCamp \href{https://drive.google.com/file/d/1xVLk15A1LBbyXH1fZnd4eBTaJ4d8FLLO/view?usp=drive_link}{}
            \item \textbf{Data Scientist Professional with R} -- DataCamp \href{https://drive.google.com/file/d/1Woe6xqxofloFZ9uM2f5VsdGxY-WQWCRI/view?usp=drive_link}{}
        \end{highlightsforbulletentries}
    \end{onecolentry}


    \section{Proyectos}

\begin{twocolentrynarrow}{\href{https://sima-d3qj5vwxtq-uc.a.run.app}{app live}}
    \textbf{SIMA -- Sistema Integral de Modelación Actuarial}
    \begin{highlights}
        \item Valuar reservas de vida exige conectar mortalidad, producto y capital en un solo flujo. SIMA lo integra: \hl{Lee-Carter}, conmutación, \hl{BEL} para tres productos y \hl{RCS} bajo \hl{LISF} con pruebas de estrés. \textit{Python, FastAPI, React, GCP.}
    \end{highlights}
\end{twocolentrynarrow}

\vspace{0.01 cm}

\begin{twocolentrynarrow}{\href{https://gonorandres.github.io/blog/data-analyst-portfolio/}{blog}}
    \textbf{Portafolio de Analista de Datos}
    \begin{highlights}
        \item Un analista se mide por el rango de preguntas que puede responder. Siete análisis end-to-end: \hl{cohortes} de e-commerce, pruebas \hl{A/B}, KPIs ejecutivos, funnels y riesgo financiero, cada uno con tablero propio. \textit{SQL, Python, Streamlit, Power BI.}
    \end{highlights}
\end{twocolentrynarrow}

\vspace{0.01 cm}

\begin{twocolentrynarrow}{\href{https://gmm-explorer.vercel.app/contexto}{app live}}
    \textbf{GMM Explorer -- Tarificador de Gastos Médicos}
    \begin{highlights}
        \item Tarificar GMM sin datos reales es especular. Procesa 5.1M siniestros de la CNSF, los clasifica en tres niveles con IA y calcula prima pura por \hl{frecuencia-severidad} ajustada por inflación médica. \textit{Python, Next.js, Vercel.}
    \end{highlights}
\end{twocolentrynarrow}

\vspace{0.01 cm}

\begin{twocolentrynarrow}{\href{https://gonorandres.github.io/blog/credit-graph-topological-risk/}{blog}}
    \textbf{CreditGraph -- Riesgo Crediticio Topológico}
    \begin{highlights}
        \item Un modelo de crédito que analiza a cada deudor por separado pierde de vista las garantías cruzadas que conectan la cartera. Modela el portafolio como red de exposiciones con \hl{Neo4j}, \hl{PySpark} en Databricks y \hl{LightGBM}+Platt. \textit{Python, Cypher, Databricks.}
    \end{highlights}
\end{twocolentrynarrow}

\vspace{0.01 cm}

\begin{twocolentrynarrow}{\href{https://gonorandres.github.io/blog/data-engineering-platform}{blog}}
    \textbf{Plataforma de Datos en GCP para Seguros}
    \begin{highlights}
        \item Un siniestro viaja lejos entre el evento y el modelo que lo tarifica. Pipeline end-to-end: \hl{Pub/Sub}+Beam, warehouse \hl{BigQuery}, \hl{Dagster}, Terraform y modelo \hl{GLM Tweedie}. \textit{Python, GCP, Docker.}
    \end{highlights}
\end{twocolentrynarrow}

\vspace{0.01 cm}

\begin{twocolentrynarrow}{\href{https://actuarial-regulation-agent-d3qj5vwxtq-uc.a.run.app/}{app live}}
    \textbf{Asistente de Regulación Actuarial (RAG)}
    \begin{highlights}
        \item LISF y CUSF suman 1,000+ artículos que Ctrl+F no resuelve. Agente que indexa ambas leyes con \hl{FTS5/BM25} y grafo de referencias y responde con cita exacta. \textit{Claude Agent SDK, FastAPI, GCP.}
    \end{highlights}
\end{twocolentrynarrow}

\end{document}
